\section{Results and Analysis}
\label{sec:results}

All numbers in this section are produced by the verified pipeline of
Section~\ref{sec:setup} under the causal information-flow protocol,
with seed 42 as the reference run and seeds $\{41,42,43\}$ as the
variance study. We release the JSON artifacts behind every table
alongside the code.

\subsection{Forecast Calibration (Stage 1)}
\label{sec:results-calibration}

\begin{figure}[!t]
\centering
\includegraphics[width=\linewidth]{audit_demand}
\caption{Data audit of the RAN demand series (1{,}778 contiguous
15-min slots): top --- the aggregate demand with train/validation/test
boundaries; bottom --- the empirical distribution, showing the
right-skew and diurnal structure that motivate quantile-based
reservations.}
\label{fig:audit}
\end{figure}

\begin{figure}[!t]
\centering
\includegraphics[width=\linewidth]{history}
\caption{Stage-1 training curves (seed 42): train and validation
pinball loss. Early stopping halts training when validation loss
stalls; the small train/val gap indicates the 52k-parameter model is
not overfitting the 1{,}175-window training segment.}
\label{fig:history}
\end{figure}

\begin{figure}[!t]
\centering
\includegraphics[width=\linewidth]{fig_quantile_band}
\caption{Stage-1 output on the test window (first 240 slots shown):
the q05--q99 forecast band and median against realized demand. The
band widens precisely where demand is volatile (morning ramp,
evening peak) and tightens overnight --- the behavior Stage~2's
buffer term converts into reservation slack.}
\label{fig:band}
\end{figure}

Fig.~\ref{fig:audit} shows the demand series and its distribution;
Fig.~\ref{fig:history} the training dynamics. Calibration is the
decision-critical property: across the 252-slot test window the
realized demand falls inside the q05--q99 band in 93.3\% of slots
(seed 42), with the three-seed mean $93.7\% \pm 3.8\%$ --- consistent
with the nominal 94\% and within the tolerance band we set a priori.
This matters operationally: the reservation in Eq.~\eqref{eq:phi} is
built from the q90 point and the q99--q50 dispersion, so a
systematically miscalibrated band would translate directly into
either violations (narrow band) or waste (inflated band).
Fig.~\ref{fig:band} shows the served forecast band against realized
demand: the band is narrow where demand is stable and wide on ramps,
which is exactly the adaptive-slack behavior the transform needs.

\subsection{Main Closed-Loop Result on Live RAN Data}
\label{sec:results-main}

\begin{table}[!t]
\caption{Reservation outcomes on the 252-slot RAN test window
(capacity 143{,}954.5; seed 42 reference run). Best value per column
in bold among deployable policies; the oracle is a cheating upper
bound and is excluded from bolding. $V$: violation rate; $U$:
utilization; $W$: waste; $E$: tracking error.}
\label{tab:main}
\centering
\small
\begin{tabular}{@{}lrrrr@{}}
\toprule
Policy & $V$ (\%) & $U$ (\%) & $W$ (\%) & $E$ (\%) \\
\midrule
Static peak          & 0.0  & 54.8 & 30.5 & 30.5 \\
Median forecast      & 48.0 & 95.1 & 1.0  & 2.0  \\
Train-only $\Phi$    & 0.0  & 54.9 & 30.3 & 30.3 \\
\textbf{UCRA (full loop)} & \textbf{0.0} & \textbf{72.1} & \textbf{10.7} & \textbf{10.7} \\
\midrule
Oracle quantile (cheating) & 16.3 & 64.9 & 19.5 & 20.1 \\
\bottomrule
\end{tabular}
\end{table}

\begin{figure}[!t]
\centering
\includegraphics[width=\linewidth]{fig_demand_vs_reservation}
\caption{The main result: UCRA's reservation tracks the diurnal
demand cycle and stays above realized demand in every one of the 252
test slots, while the static-peak baseline (dashed) pays 30.5\% of
capacity for the same zero-violation outcome.}
\label{fig:main}
\end{figure}

Table~\ref{tab:main} and Fig.~\ref{fig:main} present the central
comparison. UCRA holds \textbf{0.0\% violations} --- demand exceeds
the reservation in \emph{none} of the 252 test slots --- at
\textbf{72.1\% utilization} and \textbf{10.7\% waste}. Reading the
baselines:

\emph{Static peak} achieves the same 0.0\% violations but wastes
30.5\% of capacity --- nearly three times UCRA's waste, i.e., on this
three-week window UCRA frees up roughly 28\% of a week of capacity
that peak sizing would have idled, with no SLA cost.

\emph{Median forecast} is the cautionary tale: it has the \emph{best}
tracking error of all policies (2.0\%) --- and violates 48.0\% of
slots. A policy that reserves the median of a right-skewed demand
distribution is, by construction, insufficient half the time. This
row is the clearest demonstration of why tracking error alone must
never be the objective (Section~\ref{sec:system}).

\emph{Train-only $\Phi$} applies UCRA's own kappa-buffer rule with
training-window statistics frozen: 0.0\% violations, but 30.3\% waste
--- statistically indistinguishable from static peak. This is an
important control: the \emph{formula} alone buys nothing; the value
comes from the \emph{adaptive} Stages 1--4 updating quantiles and
dispersion per slot from recent data. The uncertainty machinery
matters only when it moves.

\emph{Oracle quantile}, which cheats by computing the rolling
quantile of realized test demand, still violates 16.3\% of slots and
achieves only 64.9\% utilization. Two lessons follow. First,
UCRA's 0.0\% at 72.1\% \emph{dominates the cheater on both axes}:
the kappa buffer plus risk feedback out-sizes a naive rolling
quantile even when that quantile sees the future it defends against.
Second, a plain q90 policy violates far more than its nominal 10\%
--- heavy-tailed 15-minute bursts push the violation probability of
a rolling historical quantile well above nominal --- confirming that
calibrated forecast-based reservation, not historical quantiles, is
the right primitive.

\subsection{The $\kappa$ Frontier}
\label{sec:results-kappa}

\begin{table}[!t]
\caption{The risk--utilization frontier traced by the risk dial
$\kappa$ (mean $\pm$ sd over seeds 41/42/43; raw $\Phi$ frontier on
the frozen forecast, no EWMA/hysteresis, $\rho = 0$).}
\label{tab:kappa}
\centering
\small
\begin{tabular}{@{}lrr@{}}
\toprule
$\kappa$ & Violation rate (\%) & Utilization (\%) \\
\midrule
0.00 & 0.40 $\pm$ 0.00 & 75.5 $\pm$ 4.1 \\
0.25 & 0.00 $\pm$ 0.00 & 71.2 $\pm$ 4.3 \\
0.50 & 0.00 $\pm$ 0.00 & 67.5 $\pm$ 4.4 \\
0.75 & 0.00 $\pm$ 0.00 & 64.3 $\pm$ 4.3 \\
1.00 & 0.00 $\pm$ 0.00 & 61.4 $\pm$ 4.3 \\
1.50 & 0.00 $\pm$ 0.00 & 56.4 $\pm$ 4.2 \\
2.00 & 0.00 $\pm$ 0.00 & 52.4 $\pm$ 4.0 \\
\bottomrule
\end{tabular}
\end{table}

\begin{figure}[!t]
\centering
\includegraphics[width=\linewidth]{fig_kappa_sweep}
\caption{The $\kappa$ sweep on the RAN test window (seed 42):
violation rate (left axis) collapses to zero by $\kappa = 0.25$,
while utilization (right axis) decays monotonically. $\kappa$ is a
pure frontier-position dial.}
\label{fig:kappa}
\end{figure}

Table~\ref{tab:kappa} and Fig.~\ref{fig:kappa} trace the frontier.
At $\kappa = 0$ (bare quantile tracking, no buffer) the raw-$\Phi$
policy violates in only 0.4\% of slots --- the calibrated q90
forecast is nearly sufficient by itself --- and each unit of
$\kappa$ buys down the remaining risk at a measured utilization
price: the knee (smallest zero-violation point) sits between
$\kappa = 0$ and $0.25$, and the deployed $\kappa = 0.5$ keeps a
zero-violation margin against seed-level variation in forecast
quality (utilization spread $\pm 4.4$ points across seeds) at a cost
of $\approx 4$ utilization points versus the knee. Beyond the knee,
$\kappa$ is a pure slack dial: violations stay at zero and utilization
decays $\approx 4$ points per 0.25 of $\kappa$. An operator with a
harder SLA can move right; a capacity-starved operator can move
left --- with the frontier quantified rather than guessed.

\subsection{Drift Response: Self-Evolution Under Surge}
\label{sec:results-drift}

\begin{table}[!t]
\caption{Drift experiment: $+35\%$ step surge on realized demand from
slot 37 of the 252-slot test window; both passes start from the same
checkpoint; the evolving pass uses the causal replay protocol.
Top: seed 42. Bottom: mean $\pm$ sd over seeds 41/42/43.}
\label{tab:drift}
\centering
\small
\begin{tabular}{@{}lrrr@{}}
\toprule
Pass & $V$ post-surge (\%) & $U$ post-surge (\%) & Updates \\
\midrule
\multicolumn{4}{@{}l}{\emph{Seed 42}} \\
Frozen (Stages 1--4) & 27.4 & 83.2 & 0 \\
Self-evolving (causal) & 12.1 & 81.1 & 3 ([96, 120, 144]) \\
\midrule
\multicolumn{4}{@{}l}{\emph{Three seeds}} \\
Frozen (Stages 1--4) & 34.3 $\pm$ 5.9 & 80.4 $\pm$ 2.8 & 0 \\
Self-evolving (causal) & 17.2 $\pm$ 4.9 & 78.2 $\pm$ 2.8 & 5.7 $\pm$ 2.3 \\
\bottomrule
\end{tabular}
\end{table}

\begin{figure}[!t]
\centering
\includegraphics[width=\linewidth]{fig_drift_demo}
\caption{Drift experiment (seed 42). Top: realized demand with the
$+35\%$ surge (dashed line), the frozen reservation (red), and the
self-evolving reservation (green) --- the green line re-sizes upward
after each fine-tune event. Bottom: rolling violation rates against
the 15\% trigger; vertical green lines mark the three causal updates
at slots 96, 120, 144.}
\label{fig:drift}
\end{figure}

Table~\ref{tab:drift} and Fig.~\ref{fig:drift} report the drift
response. The surge is a $+35\%$ step multiplier on realized demand
--- a violent regime shift of the kind a mass event produces ---
starting 15\% into the test window. The frozen forecaster keeps
reserving around the old distribution: post-surge violations reach
27.4\% (seed 42; $34.3\% \pm 5.9$ across seeds). The self-evolving
pass faces the identical surge with the identical information budget,
but its Stage~5 fires: the trigger crosses 15\% at slot 96, the
engine fine-tunes causally on replayed history, and post-surge
violations fall to 12.1\% (seed 42; $17.2\% \pm 4.9$ across seeds) ---
a \textbf{50--56\% relative reduction} --- at slightly
\emph{lower} post-surge utilization (81.1\% vs 83.2\%). The
utilization comparison matters for honesty: adaptation did not come
from blanket over-provisioning; the fine-tuned quantiles placed the
added buffer where the new regime actually demands it.

The update timeline is diagnostic of the mechanism: triggers are
evaluated every 24 slots, the rolling 96-slot violation window
crosses the 15\% trigger at slot 96, and updates fire at 96, 120,
and 144 in \emph{every} seed --- after which the window fills with
post-update, low-violation slots and the trigger silences. Seeds
41 and 43 (whose frozen baselines violate more) fire four additional
refreshes later (5.7 $\pm$ 2.3 updates total), tracking the
violation severity of each seed's own forecaster. The residual
12--17\% post-surge violations are the honest cost of a $+35\%$ step
absorbed by three gentle fine-tunes; the sensitivity grid in the
released artifacts shows more aggressive fine-tuning hyperparameters
chase the residual at the price of overfitting to spikes.

\subsection{Trigger Discipline: The Zero-Update Control}
\label{sec:results-control}

\begin{figure}[!t]
\centering
\includegraphics[width=\linewidth]{fig_self_evolution}
\caption{The clean-run control: rolling violation rate on the
untouched RAN test window never approaches the 15\% trigger; Stage~5
fires zero updates across all seeds. The insurance policy stays
silent when there is nothing to fix.}
\label{fig:control}
\end{figure}

On the \emph{clean} test window --- no injected drift --- Stage~5
never fires: zero updates in seed 42 and in every seed of the study,
with the rolling violation rate never approaching the 15\% trigger
(Fig.~\ref{fig:control}) and demand never exceeding three training
sigmas. We report this deliberately as a first-class result. A
trigger-based adaptive system that fires spuriously would erase its
own benefit through unnecessary fine-tuning churn and the risk of
adapting to noise; UCRA's dual trigger maintained perfect discipline
across 252 slots $\times$ 3 seeds of clean data while still firing
reliably and promptly under genuine drift
(Section~\ref{sec:results-drift}). This symmetric evidence ---
silence when unneeded, action when needed --- is the evaluation
standard we argue trigger-based network adaptivity should be held
to. It also bounds the claim of the paper honestly: on stationary
data, Stages 1--4 already deliver the full result, and Stage~5's
value is realized specifically under regime shift.

\emph{Runtime validation of the causal contract.} The causal
protocol is not only unit-tested; the run logs its own compliance.
In the clean RAN run, 252 windows were offered to the replay intake
and 249 released --- exactly the three windows whose final labels
fall beyond the end of the test window and would never be fully
observed by a live deployment; the label-availability delay is
recorded as $H - 1 = 3$ slots in the results artifact. Every
fine-tune batch in the drift experiment is additionally checked
against Proposition~\ref{prop:causal} by the release-time invariant,
so any future code change that reintroduces an admission race fails
the run rather than the paper.

\subsection{Seed Robustness}
\label{sec:results-seeds}

\begin{table}[!t]
\caption{Three-seed verification (seeds 41/42/43; fresh Stage-1
training per seed; identical causal evaluation). Mean $\pm$ sample
sd.}
\label{tab:seeds}
\centering
\small
\begin{tabular}{@{}lrr@{}}
\toprule
Quantity & Mean & sd \\
\midrule
q05--q99 coverage (\%) & 93.7 & 3.8 \\
UCRA violation rate (\%) & 0.0 & 0.0 \\
UCRA utilization (\%) & 67.6 & 3.4 \\
UCRA waste (\%) & 13.3 & 2.1 \\
Clean-run evolution updates & 0 & 0 \\
Demo: frozen $V$ post-surge (\%) & 34.3 & 5.9 \\
Demo: evolving $V$ post-surge (\%) & 17.2 & 4.9 \\
\bottomrule
\end{tabular}
\end{table}

Table~\ref{tab:seeds} aggregates the verification study. The
qualitative claims are seed-invariant: violations are 0.0\% in every
seed; coverage is consistent with nominal in every seed; the
drift-demo gap persists with overlapping-but-separated distributions
(frozen $34.3\% \pm 5.9$ vs evolving $17.2\% \pm 4.9$; the per-seed
pairs are $(38.1, 21.9)$, $(27.4, 12.1)$, $(37.2, 17.7)$ --- evolving
wins by 12--16 points in each seed). The quantitative spread is
concentrated exactly where one expects it: in the learned component.
Utilization varies by $\pm 3.4$ points (per-seed 63.9 / 72.1 / 66.9)
with a compensating waste spread of $\pm 2.1$ points, tracking
seed-level differences in how conservative each trained forecaster's
upper quantiles are; Stages 2--4, being deterministic arithmetic,
contribute zero additional variance. For deployment planning, the
practical reading is: pick $\kappa$ for the \emph{worst} seed's
forecast quality --- the frontier of Table~\ref{tab:kappa} already
averages over it.

\subsection{Reservation Dynamics and Churn}
\label{sec:results-dynamics}

\begin{figure}[!t]
\centering
\includegraphics[width=\linewidth]{fig_reservation_dynamics}
\caption{Reservation dynamics on the RAN test window (seed 42).
(a)~The slack $(\Rt - \D)/C$ is positive in all 252 slots, tightly
distributed (mean 10.7\%, 5th percentile 3.1\%, minimum 0.3\%) ---
the reservation is neither systematically inflated nor on the edge of
failure. (b)~Stage-4 EWMA smoothing with release hysteresis reduces
the mean per-slot reservation adjustment by 15.9\% and the maximum by
30\% relative to raw target-chasing, suppressing reservation churn.}
\label{fig:dynamics}
\end{figure}

Table~\ref{tab:main} summarizes \emph{where} the reservation sits;
Fig.~\ref{fig:dynamics} shows \emph{how} it moves. Two operational
properties follow. First, the slack distribution is tight and
strictly positive: the closest approach to a violation in 252 slots
is 0.3\% of capacity, the median headroom is 10.7\%, and the 5th
percentile is 3.1\% --- the system runs neither wastefully bloated
nor knife-edge. Neither clamp of Eq.~\eqref{eq:phi} binds on this
window (0 of 252 slots at floor or ceiling): the reservation lives
in the transform's interior, where its behavior is pure
forecast-plus-buffer arithmetic rather than clamp-driven. Second,
the actuation chain matters as much as the target: chasing the raw
$\Phi$ target slot-by-slot would adjust the reservation by up to
6.8\% of capacity in one slot, whereas the deployed
EWMA-plus-hysteresis commitment caps the same sequence at 4.7\% and
cuts mean churn by 15.9\%. In a production slice manager, every
adjustment propagates into upstream commitments and signaling; churn
suppression is not cosmetic.

\emph{Where the slack comes from.} Decomposing the committed
reservation over the run log: the base quantile term contributes on
average 91.4\% of $\Rt$ (mean base 62{,}532) and the kappa buffer
8.5\% (mean buffer 5{,}815, ranging 6.0\% to 26.2\% of the base as
the forecast's dispersion widens and narrows). The mean dispersion
is 11{,}630, i.e., 21.9\% of the mean demand --- the uncertainty the
forecaster itself reports --- and the risk modulation never
activates ($\rho_t = 0$ throughout, consistent with zero
violations). The reservation's slack is thus bought almost entirely
by the buffer term reacting to measured forecast dispersion, exactly
as designed.

\emph{Cold-start behavior.} The loop initializes its reservation from
the training window's 90\% quantile rather than from zero or from the
first forecast. On this run that seed value ($\approx$81\,203) sits
above every realized demand of the first test slots, so the system
opens with positive slack and the EWMA converges onto the forecast
trajectory within roughly a dozen slots without ever dipping below
realized demand. A cold start at the median forecast, by contrast,
opens the evaluation window with immediate violations and pollutes
the drift monitor's early violation statistics --- an avoidable
configuration error that the released configuration forecloses.

\subsection{Drift-Severity Grid and Update Sensitivity}
\label{sec:results-grid}

\begin{table}[!t]
\caption{Drift-severity grid (seed 42; surge applied from slot 37;
causal protocol throughout). Relative reduction $=$ frozen minus
evolving, normalized by frozen. The trigger threshold creates a
deliberate no-adaptation regime at mild surges: when post-surge
violations stay below the 15\% trigger, Stage~5 correctly stays
silent.}
\label{tab:severity}
\centering
\small
\begin{tabular}{@{}lrrrr@{}}
\toprule
Surge & Frozen $V$ & Evolving $V$ & Rel.\ red. & Updates \\
\midrule
$+15\%$ & 4.2\%  & 4.2\%  & 0.0\%  & 0 \\
$+25\%$ & 12.1\% & 12.1\% & 0.0\%  & 0 \\
$+35\%$ & 27.4\% & 12.1\% & 55.9\% & 3 ([96, 120, 144]) \\
$+50\%$ & 53.5\% & 26.5\% & 50.4\% & 8 ([72, \dots, 240]) \\
\bottomrule
\end{tabular}
\end{table}

\begin{table}[!t]
\caption{Fine-tune sensitivity at the reference surge ($+35\%$, seed
42): increasing fine-tune epochs from 3 to 10 narrows the residual
post-surge violations (12.1\% $\rightarrow$ 9.8\%) with identical
trigger points, confirming the behavior is a property of the
trigger-plus-replay design rather than of one learning rate.}
\label{tab:ftepochs}
\centering
\small
\begin{tabular}{@{}lrrr@{}}
\toprule
Fine-tune epochs & Evolving $V$ (\%) & Rel.\ red. & Updates \\
\midrule
3 (default) & 12.1 & 55.9\% & [96, 120, 144] \\
10          & 9.8  & 64.4\% & [96, 120, 144] \\
\bottomrule
\end{tabular}
\end{table}

Table~\ref{tab:severity} sweeps the surge magnitude and exposes the
policy structure honestly. At $+15\%$ and $+25\%$, the frozen
pipeline's post-surge violations (4.2\% and 12.1\%) remain below the
15\% trigger, so Stage~5 fires zero updates and the evolving pass is
identical to the frozen one --- by design, the system tolerates
mild surges with its existing slack and does not spend updates on
noise. At $+35\%$ the trigger engages and adaptation halves
violations; at $+50\%$ the trigger fires earlier (slot 72) and eight
times, and while absolute violations remain high (26.5\%), the
relative reduction persists (50.4\%) --- a violent step cannot be
fully absorbed by three-to-eight gentle fine-tunes, and we report
the residual rather than tuning it away. Table~\ref{tab:ftepochs}
addresses the natural reviewer question of whether the adaptation
depends on fine-tuning intensity: with 10 epochs instead of 3, the
same trigger points produce a deeper recovery (9.8\% vs 12.1\%),
i.e., the mechanism is robust and its intensity is itself a tunable
knob.


\subsection{Slice-Level Validation on CTTC}
\label{sec:results-cttc}

\begin{table}[!t]
\caption{CTTC slice-level comparison (229 retained snapshots from
250 parsed; causal serving loop; train quantiles from the 70\%
train pool). $V$: violation rate; $U$: utilization; $O$: over- and
under-provisioning measured as mean unused share of the reservation.}
\label{tab:cttc}
\centering
\small
\setlength{\tabcolsep}{3.5pt}
\begin{tabular}{@{}llrrr@{}}
\toprule
Slice & Policy & $V$ (\%) & $U$ (\%) & $O$ (\%) \\
\midrule
\multirow{4}{*}{URLLC}
 & Operator ($\delta$-sizing) & 29.2 & 54.9 & 45.1 \\
 & Static $q_{90}$           & 12.1 & 43.3 & 56.7 \\
 & Train-only $\Phi$         & 1.2  & 18.6 & 81.4 \\
 & \textbf{UCRA $\Phi$ (causal)} & \textbf{1.2} & 18.6 & 81.4 \\
\midrule
\multirow{4}{*}{eMBB}
 & Operator ($\delta$-sizing) & 43.7 & 71.5 & 28.5 \\
 & Static $q_{90}$           & 4.4  & 40.2 & 59.8 \\
 & Train-only $\Phi$         & 1.9  & 21.4 & 78.6 \\
 & \textbf{UCRA $\Phi$ (causal)} & \textbf{1.9} & 21.3 & 78.7 \\
\midrule
\multirow{4}{*}{mMTC}
 & Operator ($\delta$-sizing) & 5.8  & 19.9 & 80.1 \\
 & Static $q_{90}$           & 17.6 & 42.2 & 57.8 \\
 & Train-only $\Phi$         & 2.6  & 18.5 & 81.5 \\
 & \textbf{UCRA $\Phi$ (causal)} & \textbf{2.6} & 18.4 & 81.6 \\
\bottomrule
\end{tabular}
\end{table}

\begin{figure}[!t]
\centering
\includegraphics[width=\linewidth]{fig_cttc_slices}
\caption{CTTC slice-level outcomes by policy: violation rate (left)
and over-provisioning (right). The operator's default $\delta$-sizing
violates heavily for URLLC and eMBB; static $q_{90}$ is
inconsistent across slice types; the kappa-buffered transform holds
1.2--2.7\% violations everywhere at a measured over-provisioning
price.}
\label{fig:cttc}
\end{figure}

Table~\ref{tab:cttc} and Fig.~\ref{fig:cttc} validate the transform
on an independent workload with per-slice ground truth. Three
findings. First, the operator's default $\delta$-sizing violates
substantially on the two performance-critical slice types --- 29.2\%
of URLLC and 43.7\% of eMBB snapshots see offered load beyond the
reservation --- while being profligate on mMTC (80\% unused). Second,
a bare empirical $q_{90}$ (the $\kappa = 0$ point of $\Phi$) is
unreliable across slice types: it under-protects URLLC (12.1\%
violations) and mMTC (17.6\%) --- heavy-tailed offered load again
pushing realized violations far above the nominal 10\% --- while the
kappa-buffered transform holds \textbf{1.2\% (URLLC), 1.9\% (eMBB),
2.6\% (mMTC)}, at the price of high measured over-provisioning
($\approx$80\%), the same frontier shape as the RAN side. Third, the
honest negative result: on this dataset UCRA $\Phi$ is
indistinguishable from train-only $\Phi$ --- the online risk
indicator $\rho_t$ stays near zero because the snapshot stream
exhibits no sustained violation runs for the feedback to amplify.
We therefore attribute the CTTC gains to the kappa buffer acting on
train-quantile dispersion, not to risk modulation, and note that the
RAN-side experiments are where the adaptive machinery demonstrably
earns its complexity.

\emph{A note on QoS outcomes.} We also examined whether
operator-reservation violations map to worse observed QoS at the
snapshot level. They do not, cleanly: mean packet-drop ratios of
violated versus satisfied snapshots are statistically
indistinguishable for URLLC and eMBB, and lower for violated mMTC
snapshots (1.0\% vs 2.4\%), while the delay fields in this dataset
are log-scaled and similarly flat. The most plausible reading is
that steady-state snapshots average over the transient congestion
that a reservation violation represents, so the snapshot-level QoS
metrics are too coarse an instrument to grade violation severity.
We therefore treat the violation rate itself --- the quantity the
reservation policy directly controls --- as the primary outcome on
this dataset, and flag snapshot-QoS correlation as an open
measurement question rather than asserting it.
